(e) From (a) it is easy to see that
\[
[x^\alpha y^\beta]s_\lambda(x/y)=\langle h_\alpha e_\beta,s_\lambda\rangle.
\]
This corresponds to building up the shape $\lambda$ by first adding horizontal strips of sizes $\alpha_1,\alpha_2,\ldots$ and then vertical strips of sizes $\beta_1,\beta_2,\ldots$\,. But multiplication of symmetric functions is \emph{commutative}, so we could first add a horizontal strip $H_1$ of size $\alpha_1$, then a vertical strip $V_1$ of size $\beta_1$, then a horizontal strip $H_2$ of size $\alpha_2$, then a vertical strip $V_2$ of size $\beta_2$, etc. To get $s_\lambda(x/x)$, fill in each square of $H_i\cup V_i$ with $i$, getting an array $T$, and sum all $x^T$'s. Now the array $T$ is precisely a supertableau. How many times does a particular supertableau $T$ get counted? For each component $C$ of $H_i\cup V_i$, there are exactly two ways to decompose $C$ into a horizontal strip $H$ followed by a vertical strip $V$\,. (The last square in the top row can be part of $H$ or $V$\,, but there are no other choices.) Hence the number of times $T$ gets counted is $2^{c(T)}$, completing the proof.
